\section{Related Works}
\label{sec:formatting}

\paragraph{Audio-driven facial animation}
The goal of audio-driven facial animation methods is to generate high-quality talking head videos that preserve the identity of input faces while ensuring accurate lip movements synchronized with the input audio. Early GAN-based methods~\cite{DBLP:conf/bmvc/VougioukasPP18,Vougioukas, zhou2019talking, chung2017you} mostly focused on animating the speaker's facial expressions by introducing temporal constraints and expert discriminators to improve lip-sync accuracy. Later, several works~\cite{chen2020talking, zhang2022sadtalker, zhou2021pose} built on these approaches by incorporating head pose modelling to generate more realistic animations, but were prone to producing artifacts and unnatural motion.

Diffusion models~\cite{ho2020ddpm, rombach2022highresolutionimagesynthesislatent} have emerged as an alternative to GANs for audio-driven facial animation, demonstrating improved temporal consistency and video quality~\cite{du2023dae, xu2024vasa}. Several methods~\cite{stypulkowski2023, xu2024hallohierarchicalaudiodrivenvisual, wang2024V-Express, chen2024echomimiclifelikeaudiodrivenportrait} leverage video diffusion models~\cite{ho2022video, blattmann2023stablevideodiffusionscaling} for temporally consistent motion. Additionally, some works propose to condition the generation process on facial landmarks~\cite{wei2024aniportraitaudiodrivensynthesisphotorealistic} or 3D meshes \cite{zhang2023dream}. However, these approaches often produce non-realistic facial motion. Finally, a recent line of works~\cite{wei2024aniportraitaudiodrivensynthesisphotorealistic, wang2024V-Express, chen2024echomimiclifelikeaudiodrivenportrait} leverage ReferenceNet~\cite{animateanyone} to improve identity reconstruction, though at the cost of increased computational complexity. 

Recently, KeyFace~\cite{keyface} introduced a keyframe-based approach that predicts key poses and interpolates between them, enhancing identity preservation and temporal consistency. We follow their strategy, tailoring it to lip synchronization to ensure temporally consistent lip-sync animations that preserve the original identity without visible inpainting borders or expression leakage from the input video.

\begin{figure*}
  \centering
  \includegraphics[width=\textwidth]{Images/drawing.pdf}
    \caption{\textbf{Overview of the KeySync framework.} KeySync consists of two stages, both of which involve generating video using latent diffusion conditioned on an input video and audio, differing only in the reference frames selection, as described in (b). During keyframe generation, the model receives an identity frame  $x_{\text{id}}$, which is repeated and concatenated with the noised video input. During interpolation, the model is conditioned on two successive keyframes $z_{i}$ and $z_{i+1}$, along with intermediate learnable embeddings $z_m$. Both stages integrate audio embeddings $a$ from HuBERT~\cite{hubert}. In (c), we illustrate our occlusion handling pipeline, which we apply during inference.}
    \label{fig:architecture}
\end{figure*}

\paragraph{Audio-driven lip synchronization}
Lip synchronization methods focus on adjusting mouth movements to match the input audio while preserving other facial attributes, \ie, the head pose and upper face expressions. The first notable work, Wav2Lip \cite{Prajwal}, uses GANs to generate a sequence of frames from input video frames with the lower part of the face masked. To improve lip synchronization, Wav2Lip~\cite{Prajwal} leverages a pre-trained lip-sync expert model. In order to enhance realism and identity generalization, StyleSync~\cite{guan2023stylesync} and StyleLipSync~\cite{StyleLipSync} introduce StyleGAN2-based~\cite{karras2020analyzing} architectures, while DINet~\cite{DINet} performs spatial deformation on feature maps to improve visual quality. Finally, TalkLip~\cite{Zhong_2023_CVPR} proposes using contrastive learning based on a pre-trained lip-sync expert to enhance the quality and accuracy of generated lip region.

Recently, diffusion-based methods for lip synchronization have been introduced~\cite{diff2lip, DiffDub, michaldub}. Nevertheless, expression leakage from the input video, especially in cross-driving scenarios, remains an open issue. Several approaches attempt to mitigate this using different masking strategies~\cite{StyleLipSync, DINet, zhang_musetalk_2024, VideoReTalking}, but no consensus exists regarding the optimal masking method. Another challenge is temporal consistency, as many methods~\cite{maketalk, DiffDub, stylepres} operate on a frame-by-frame basis without explicit sequence modeling, leading to discontinuities. Some models are conditioned on past frames~\cite{michaldub}, but consequently suffer from cumulative error propagation, while others propose the use of pre-trained perceptual models~\cite{latentsync} or sequence discriminators~\cite{diff2lip} to enforce coherence, but are generally not sufficient. Finally, occlusion handling remains an open challenge, as most models assume clear visibility of the mouth, failing in real-world settings where occlusions from hands, objects, or motion blur occur.

Inspired by KeyFace~\cite{keyface}, we propose a two-stage lip synchronization framework ensuring robust, leakage-free cross-driving performance. Furthermore, our post-training occlusion-handling strategy further enhances robustness, rendering our model suitable for real-world applications.
